\documentclass[10pt,reqno]{amsart}
\usepackage[letterpaper,hmargin=1.04in,vmargin=0.88in]{geometry}
\usepackage{amsmath,amssymb}
\usepackage[hidelinks]{hyperref}
\hypersetup{
  pdftitle={Nonconvexity of the product belief-propagation image for the antiferromagnetic Potts model},
  pdfauthor={Kazeshima},
  pdfkeywords={antiferromagnetic Potts model, belief propagation, product measures, convexity}
}
\newtheorem{theorem}{Theorem}
\newcommand{\E}{\mathbb E}
\newcommand{\Fprod}{\mathcal F^{\mathrm{prod}}}

\title[Nonconvexity of the product Potts BP image]{Nonconvexity of the product belief-propagation image\\
for the antiferromagnetic Potts model}
\author{Kazeshima}
\date{November 25, 2026}

\begin{document}
\begin{abstract}
For the four-state antiferromagnetic Potts model with two neighbours and
diagonal interaction $1/2$, we give two product inputs whose belief-propagation
outputs have an unattainable midpoint. This answers Question~4 of Liu, Mani
and Pernice in the negative. More generally, the product image is nonconvex
for every $d\ge2$, $q\ge d+2$ and $0\le\beta<1$. Attaining the largest
possible output probability for one colour forces the neighbours to have
pairwise disjoint supports.
\end{abstract}
\maketitle

\section{A missing midpoint in four colours}

Take two independent neighbours with respective colour distributions
$p,r\in\Delta_4$, where $\Delta_q$ is the closed probability simplex on
$[q]=\{1,\ldots,q\}$. A neighbour contributes a factor $1/2$ when its
colour agrees with the root colour and a factor $1$ otherwise. The
unnormalized weight for root colour $i$ and the normalizing sum are therefore
\begin{equation}\label{eq:four}
 g_i=(1-p_i/2)(1-r_i/2),\qquad
 Z=\sum_{i=1}^4g_i=3+\tfrac14\sum_{i=1}^4p_i r_i .
\end{equation}
Write $F(p,r)=(g_i/Z)_{i=1}^4$ and let $\delta_i$ denote the point mass at $i$.
Two choices of messages give
\[
 u=F(\delta_1,\delta_2)=\tfrac16(1,1,2,2),\qquad
 v=F\bigl(\delta_3,(\delta_1+\delta_2)/2\bigr)=\tfrac1{12}(3,3,2,4).
\]
Their midpoint is $w=(5,5,6,8)/24$. We claim that $w$ is unattainable.

Indeed, $g_4\le1$ and $Z\ge3$, so $F_4\le1/3$. Since $w_4=1/3$, any
preimage of $w$ must have both $g_4=1$ and $Z=3$. The first equality
forces $p_4=r_4=0$. The second forces $p_i r_i=0$ for every $i$, by
\eqref{eq:four}. Thus on the three remaining colours
$g_i=1-(p_i+r_i)/2=3w_i$, giving
\[
 (p_1+r_1,\ p_2+r_2,\ p_3+r_3)=(3/4,3/4,1/2).
\]
Each positive mass belongs entirely to one neighbour. To give each
neighbour total mass $1$, the three masses must split into two nonempty
groups. One group would be a singleton, yet none of the masses is $1$.

Liu, Mani and Pernice ask whether the product belief-propagation image is
convex for $q\ge3$ and $\max\{0,1-q/(d+1)\}\le\beta\le1$
\cite[Definition~2 and Question~4]{LMP}. The same question appears in
Mani's thesis \cite[Question~3.9.18]{Mani}. The admissible interval for
$(q,d)=(4,2)$ is $[0,1]$, so the two messages above give a counterexample
at an interior value of the interaction parameter.

\section{The general obstruction}

Let $A_{ib}=\beta$ for $b=i$ and $A_{ib}=1$ otherwise, and put $t=1-\beta$.
For a joint law $\mu$ of the neighbour configuration
$\tau=(\tau_1,\ldots,\tau_d)\in[q]^d$, the map in
\cite[eqs.~(1)--(2)]{LMP} is $F_i(\mu)=G_i(\mu)/Z(\mu)$, where
\[
 G_i(\mu)=\E_\mu\prod_{j=1}^d A_{i,\tau_j},
 \qquad Z(\mu)=\sum_{b=1}^q G_b(\mu).
\]
If $\mu=\bigotimes_{j=1}^d p_j$ with arbitrary $p_j\in\Delta_q$, independence
gives $G_i(\mu)=\prod_j(1-tp_j(i))$. Denote the set of all such normalized
outputs by $\Fprod_{A,d}$.

\begin{theorem}\label{thm:main}
For all integers $d\ge2$ and $q\ge d+2$ and all real $0\le\beta<1$,
the set $\Fprod_{A,d}$ is nonconvex.
\end{theorem}

\begin{proof}
Write $Z_0=q-dt\ge q-d\ge2$. For a configuration $\tau$ let $N_i$ count
its occurrences of colour $i$. Then $G_i(\mu)=\E_\mu[\beta^{N_i}]$, with
$\beta^0=1$ also at $\beta=0$. For integers $n\ge1$,
\[
 1-\beta^n=t\sum_{k=0}^{n-1}\beta^k\le nt .
\]
Thus $\beta^n\ge1-nt$ for every $n\ge0$, with equality precisely for
$n=0,1$. Summing over colours and using $\sum_iN_i=d$ gives
\begin{equation}\label{eq:partition}
 Z(\mu)\ge Z_0 .
\end{equation}
Equality holds precisely when the neighbours have pairwise distinct colours
almost surely. For a product law this is equivalent to pairwise disjoint
supports of the $p_j$. A shared colour $i$ would be drawn by two specified
neighbours with probability $p_j(i)p_k(i)>0$. At equality, every colour
therefore satisfies
\begin{equation}\label{eq:linear}
 G_i(\mu)=1-t\sum_{j=1}^d p_j(i).
\end{equation}

Fix a set $S$ of $d+1$ colours. For each $a\in S$, take a deterministic
input which uses each colour in $S\setminus\{a\}$ once. Its weights are
$\beta$ on the occupied colours and $1$ elsewhere, with common sum $Z_0$.
The arithmetic mean $\bar w$ of these $d+1$ outputs satisfies
\[
 Z_0\bar w_i=
 \begin{cases}
  1-td/(d+1),& i\in S,\\
  1,& i\notin S .
 \end{cases}
\]
Choose $b\notin S$, which is possible since $q\ge d+2$. A product preimage
of $\bar w$ would have $G_b/Z=1/Z_0$. As $G_b\le1$, \eqref{eq:partition}
forces $Z=Z_0$. Equation~\eqref{eq:linear} now gives
$\sum_jp_j(i)=d/(d+1)$ for each $i\in S$. Exactly one neighbour must carry
each such mass. Since $d/(d+1)>1/2$, a neighbour can carry at most one of
them. There are $d+1$ masses and only $d$ neighbours. Hence this finite
convex combination lies outside $\Fprod_{A,d}$.
\end{proof}

\section{Scope and verification}

Liu, Mani and Pernice use an extremal coordinate for an auxiliary matrix
\cite[Lemma~2.2 and Remark~3]{LMP}. They identify the minimizers of a
special-colour coordinate. For the standard antiferromagnetic matrix,
equality at an upper face gives the disjoint-support condition. We found
no earlier resolution of Question~4 in the public sources examined through
November~25, 2026.

The result includes proper colourings at $\beta=0$ and persists for strictly
positive input marginals. By continuity and \eqref{eq:partition}, the image
of the closed product of simplices is compact and has positive distance
from $\bar w$. Mixing each input marginal with a sufficiently small
proportion of the uniform law therefore keeps the mean output outside the
image. The range $q\le d+1$ remains beyond this argument. Theorem~1.4
of \cite{LMP} treats $d\ge2q$ under lower tail bounds. Jerrum's conjecture
on sampling graph colourings \cite[Conjecture~1]{LMP} remains a separate
algorithmic question.

The \href{https://github.com/Kazeshima/solved-problems/tree/main/problems/potts-bp-nonconvexity}{companion Lean project}
formalizes the four-colour example and Theorem~\ref{thm:main} over the
reals, with a proof that the finite expectation for a product law equals
the factored expression above. Lean~4.19.0 and Mathlib~v4.19.0 check the
theorems. Their transitive axiom lists contain only propositional
extensionality, classical choice and quotient soundness. The repository
provides the kernel replay and reproduction commands. Generative AI
assisted in preparing the manuscript and the formalization.

\begin{thebibliography}{9}
\bibitem{LMP}
K. Liu, N. Mani and F. Pernice,
\emph{Counterexamples to a Weitz-Style Reduction for Multispin Systems},
\href{https://arxiv.org/abs/2411.06541v2}{arXiv:2411.06541v2}, 2025.

\bibitem{Mani}
N. Mani, \emph{A Probabilistic Perspective on Graph Coloring},
Ph.D. thesis, Massachusetts Institute of Technology, 2025,
\href{https://dspace.mit.edu/handle/1721.1/159909}{hdl:1721.1/159909}.
\end{thebibliography}
\end{document}
